\chapter{Нечётный фон оператора Дирака и кубическая проекция}
\label{ch:v8-baryon-noncentral-odd-dirac-background-projection-no-go-gate}

Предыдущая глава нашла единую положительную форму, но её центральный фон
оказался чётным. Проверим последнюю локальную возможность: заменить его
произвольным нецентральным нечётным самосопряжённым фоном и потребовать
проекцию на весь канонический тензор $d_{abc}$.

\section{Градуировка полного кадра}

На разложении $\mathbb C^{21}=\mathbb C^{11}\oplus\mathbb C^{10}$ положим
\begin{equation}
 \Gamma=I_{11}\oplus(-I_{10}),\qquad
 \Gamma\widehat F_a\Gamma=(-1)^{p_a}\widehat F_a,
 \quad
 p_a=\begin{cases}1,&a<30,\\0,&a\ge30.\end{cases}
 \label{eq:v8-baryon-full-frame-finite-parity}
\end{equation}
Тридцать направлений переноса нечётны, двенадцать калибровочных направлений
чётны. Любой допустимый конечный фон оператора Дирака имеет вид
\begin{equation}
 D=\begin{pmatrix}0&B^*\\B&0\end{pmatrix},
 \qquad \Gamma D\Gamma=-D.
 \label{eq:v8-baryon-general-odd-dirac-background}
\end{equation}

\section{Дополнительность кубических каналов}

Канонический следовой тензор имеет уже доказанную опору
\begin{equation}
 \operatorname{supp}d=140\,TTG+28\,GGG.
 \label{eq:v8-baryon-canonical-cubic-even-support}
\end{equation}
Каждое ненулевое произведение в этой опоре содержит чётное число нечётных
конечных множителей и потому само чётно.

Кубический тензор, создаваемый фоном $D$, определяется поляризацией
\begin{equation}
 c^{D}_{abc}
 =\operatorname{Sym}_{abc}\Tr
 \left(D\widehat F_a\widehat F_b\widehat F_c\right).
 \label{eq:v8-baryon-odd-background-cubic-tensor}
\end{equation}
Для любого чётного произведения $S_{\rm even}$ имеем
\begin{equation}
 \Tr(D_{\rm odd}S_{\rm even})=0.
 \label{eq:v8-baryon-odd-even-trace-orthogonality}
\end{equation}
Следовательно, $c^D_{abc}$ тождественно равен нулю на всей опоре $d_{abc}$
при любом выборе матрицы $B$.

На существующем физическом инцидентном фоне точный полный перебор даёт
\begin{equation}
 \operatorname{supp}c^{D_0}=130\,TTT+35\,TGG,
 \qquad
 c^{D_0}_{TTG}=c^{D_0}_{GGG}=0.
 \label{eq:v8-baryon-incidence-odd-background-support}
\end{equation}
Две опоры не просто различаются: они лежат в противоположных конечных
чётностях.

\begin{theorem}[Градуировочный запрет нечётного фонового спасения]
Никакой нечётный самосопряжённый конечный фон не может породить ненулевый
скалярный множитель всего тензора $d_{abc}$ через кубическую часть
$\Tr((D+Z)^2-D^2)^2$. Канонический тензор поддержан на чётных произведениях,
а фоновая кубическая форма --- только на нечётных.
\end{theorem}

\section{Что действительно даёт полная суперкривизна}

В полной пространственно-временной градуировке калибровочная связность
чётна как матрица, но является одноформой, тогда как поле переноса нечётно
как матрица и является нуль-формой:
\begin{equation}
 \mathbb A=d+A^{(1)}_{\rm even}+\Phi^{(0)}_{\rm odd},
 \qquad |\mathbb A|_{\rm total}=1.
 \label{eq:v8-baryon-total-odd-superconnection}
\end{equation}
Её кривизна имеет стандартное разложение
\begin{equation}
 \mathbb F=F_A+D_A\Phi+\Phi^2.
 \label{eq:v8-baryon-full-supercurvature-decomposition}
\end{equation}
Кубические каналы $GGG$ и $TTG$ в норме кривизны возникают соответственно
из членов типа
\begin{equation}
 2\langle dA,A^2\rangle,\qquad
 2\langle dA,\Phi^2\rangle
 +2\langle d\Phi,[A,\Phi]\rangle.
 \label{eq:v8-baryon-derivative-cubic-supercurvature-vertices}
\end{equation}
Они содержат пространственно-временную производную. На постоянном
нулеимпульсном срезе
\begin{equation}
 dA=d\Phi=0
 \quad\Longrightarrow\quad
 S^{(3)}_{\rm supercurvature}=0.
 \label{eq:v8-baryon-zero-momentum-supercurvature-cubic-vanishing}
\end{equation}
Поэтому совпадение секторной опоры не превращает производную вершину в
статический трёхкопийный оператор $\lambda_3W_3$.

\begin{nogocriterion}[Граница локальной суперкривизностной ветви]
После отказа центрального чётного фона и общего нечётного фонового класса
локальная беспроизводная суперкривизностная ветвь закрыта. Продолжение
допустимо только через явное производное или нелокальное шеститочечное ядро
с проверкой нулеимпульсного предела; новый постоянный коэффициент при $W_3$
вводить запрещено.
\end{nogocriterion}

\begin{protocol}
\begin{enumerate}
 \item поле вычислений: \textbf{$\mathbb Q$};
 \item числа с плавающей точкой: \textbf{отсутствуют};
 \item нечётных направлений кадра: \textbf{$30$};
 \item чётных направлений кадра: \textbf{$12$};
 \item опора $d_{abc}$: \textbf{$140\,TTG+28\,GGG$};
 \item опора физического нечётного фона:
       \textbf{$130\,TTT+35\,TGG$};
 \item пересечение опор: \textbf{пусто};
 \item нецентральный нечётный фон воспроизводит $d$: \textbf{нет};
 \item стандартные кубические вершины суперкривизны:
       \textbf{производные};
 \item постоянный нулеимпульсный кубический член: \textbf{$0$};
 \item статический коэффициент $\lambda_3$: \textbf{не выведен}.
\end{enumerate}
\end{protocol}

Машинный сертификат сохранён в
\path{s2t_v8_baryon_noncentral_odd_dirac_background_projection_no_go_gate_results.json}.

\begin{thebibliography}{9}
\bibitem{v8-baryon-odd-background-tolksdorf}
J.~Tolksdorf,
\emph{The Einstein--Hilbert--Yang--Mills--Higgs Action and the Dirac--Yukawa Operator},
arXiv:hep-th/9612149.
\bibitem{v8-baryon-odd-background-roepstorff}
G.~Roepstorff, Ch.~Vehns,
\emph{Generalized Dirac Operators and Superconnections},
arXiv:math-ph/9911006.
\end{thebibliography}